\section{Introducción}
La medición de temperatura constituye una parte fundamental de la 
instrumentación biomédica, ya que permite obtener información sobre el estado
 y funcionamiento de distintos sistemas y dispositivos utilizados en el ámbito de la salud. 
 Para realizar esta medición se emplean sensores de temperatura, dispositivos capaces de 
 transformar una variación de temperatura en una variación de alguna magnitud eléctrica, como 
 corriente o tensión, que pueda ser posteriormente adquirida y procesada por un 
 sistema electrónico. Entre los diferentes tipos de sensores existentes se encuentran los termopares, 
los RTD (Resistance Temperature Detectors) y los termistores, cada uno con sus características 
particulares de sensibilidad, linealidad y rango de medición.

En este trabajo se estudia un termistor NTC, cuyo funcionamiento se basa 
en la variación de su resistencia eléctrica con la temperatura. Al tratarse de un
 dispositivo con coeficiente de temperatura negativo, su resistencia disminuye a 
medida que aumenta la temperatura. Esta relación, además, es de carácter no lineal,
 por lo que la resistencia medida debe relacionarse con la temperatura mediante un
  modelo matemático adecuado. Para ello, se utiliza la ecuación de Steinhart-Hart,
que permite obtener la temperatura del termistor a partir de su resistencia eléctrica.

La práctica realizada consistió específicamente en la implementación de un circuito
 de medición de temperatura utilizando el termistor NTC, una resistencia de 10 k$\Omega$
y un Arduino UNO, montados sobre una protoboard. La tensión generada por el divisor resistivo 
formado por el NTC y la resistencia fija es adquirida por el Arduino
 y convertida en un valor digital mediante su conversor analógico-digital (ADC). 
A partir de esta medición, el programa permite determinar la resistencia del termistor y, 
utilizando la ecuación de Steinhart-Hart, estimar la temperatura correspondiente.

La finalidad de este trabajo es evaluar el funcionamiento de los sensores de temperatura utilizados,
analizando cómo varía su resistencia eléctrica frente a diferentes temperaturas. 
Para ello, los sensores fueron sometidos a tres condiciones de temperatura: 
temperatura ambiente, agua caliente con una temperatura aproximada de 80 °C y agua con hielo, cercana a 0 °C. 
A partir de las mediciones realizadas en cada condición, se busca observar la relación entre la temperatura
 y la resistencia de los sensores, y analizar su comportamiento frente a los cambios de temperatura. 
 De esta manera, la experiencia permite comprender de forma práctica el principio de funcionamiento
  de este tipo de sensores y evaluar su respuesta ante distintas condiciones térmicas.

En este sentido, resulta relevante destacar que, si bien los termistores NTC no constituyen
elementos lineales como los sensores RTD, su linealización mediante modelos como el de
Steinhart-Hart resulta comparativamente sencilla de implementar por hardware o software,
lo cual justifica su amplia utilización en equipamiento biomédico tal como monitores
multiparam\'etricos, incubadoras y cat\'eteres de termodiluci\'on. Sin embargo, dado que
la precisi\'on de la temperatura estimada depende directamente de la exactitud de los
coeficientes empleados en dicho modelo, cobra especial importancia analizar en qu\'e
medida la utilizaci\'on de coeficientes calculados de manera experimental, frente a
los provistos por el fabricante, puede afectar la fiabilidad de las mediciones obtenidas.